\documentclass[11pt]{article}
\usepackage[margin=0.9in]{geometry}
\usepackage{amsmath,amssymb,amsthm,mathtools}
\newtheorem{proposition}{Proposition}
\newtheorem{theorem}[proposition]{Theorem}
\newcommand{\R}{\mathbb R}
\title{From the canonical source cone to the full quantum polar}
\author{}\date{}
\begin{document}\maketitle\vspace{-2em}

Write an expectation vector and its homogenization as
\[
 q=(A_0,A_1,B_0,B_1,E_{00},E_{01},E_{10},E_{11}),\qquad z=(1,q).
\]
Let $P=\{f\in\R^9:f^Tz\ge0\text{ for every }q\in\mathcal Q_{222}\}$.
The exact geometric reduction below applies to any canonical SDP cone.
An exact canonical cone with a strict feasible lift gives the full
eight-coordinate quantum set through eight adjoint copies. The
order-$33$ specialization at the end is only an outer approximation.

\paragraph{The canonical map.}
Use $H=1-(E_{00}-E_{01}+E_{10}+E_{11})/2$ and
\[
 4p(a,b\mid x,y)=1+aA_x+bB_y+abE_{xy}.
\]
In the source order $c=(d,\lambda,n,m,C,A,D,\alpha,\beta)$, the event
labels are
\[
\begin{array}{c|cccccccc}
 &\lambda&n&m&C&A&D&\alpha&\beta\\\hline
 p_i&++|01&+-|00&+-|11&-+|10&-+|00&-+|11&--|01&+-|10
\end{array}
\]
and $f_c=dH+4\sum_i c_i p_i$. In the expectation order above, with
the constant first, the nine coefficients $f=Mc$ are exactly
\begin{align}
 f_0&=d+\lambda+n+m+C+A+D+\alpha+\beta,\notag\\
 f_{A_0}&=\lambda+n-A-\alpha,&
 f_{A_1}&=m-C-D+\beta,\notag\\
 f_{B_0}&=-n+C+A-\beta,&
 f_{B_1}&=\lambda-m+D-\alpha,\notag\\
 f_{E_{00}}&=-d/2-n-A,& f_{E_{01}}&=d/2+\lambda+\alpha,\notag\\
 f_{E_{10}}&=-d/2-C-\beta,& f_{E_{11}}&=-d/2-m-D.
                                                               \label{eq:map}
\end{align}
The map is invertible, with $\det M=-32$. Its inverse has a direct
evaluation formula. Let
$r=(0,0,0,0,1,-1,1,1)$ be the canonical PR point and
$h=\left(\begin{smallmatrix}1&-1\\1&1\end{smallmatrix}\right)$.
For the losing event $(a,b\mid x,y)$, its deterministic neighbour is
specified explicitly by
\[
 A_x=a,\quad B_y=b,\quad
 B_{1-y}=h_{x,1-y}a,\quad A_{1-x}=h_{1-x,y}b,
 \qquad E_{ij}=A_iB_j.
\]
These neighbours obey $H(v_i)=0$ and $p_j(v_i)=\delta_{ij}$;
every event $p_i$ vanishes at $r$. Hence
\begin{equation}
 d=-f(r),\qquad c_i=f(v_i)/4.                            \label{eq:inverse}
\end{equation}
Equivalently, if $V$ has rows $(1,r)$ and $(1,v_i)$, then
$VM=\operatorname{diag}(-1,4I_8)$ and
$M^{-1}=\operatorname{diag}(-1,\tfrac14 I_8)V$.

\paragraph{Eight exact relabelings.}
For $(r_0,r_1,s_0)\in\{\pm1\}^3$, flip the outputs by
$A_0\mapsto r_0A_0$, $A_1\mapsto r_1A_1$, $B_0\mapsto s_0B_0$,
leaving $B_1$ unchanged. Both homogeneous expectations and coefficient
vectors transform by
\begin{equation}
 R_g=\operatorname{diag}(1,r_0,r_1,s_0,1,r_0s_0,r_0,r_1s_0,r_1).
                                                               \label{eq:flips}
\end{equation}
These are involutions preserving $\mathcal Q_{222}$. They send the
canonical PR correlation pattern to
\[
                  (r_0s_0,-r_0,r_1s_0,r_1),
\]
giving all eight patterns with product $-1$.
Thus the complete rational maps for chart $g$ are
\begin{equation}
 M_g=R_gM,\qquad c=M^{-1}R_gf,\qquad
 M_g^Tz=M^TR_gz.                                        \label{eq:chart-map}
\end{equation}
No input interchange or party swap is required.

\begin{proposition}\label{prop:polar-cover}
Let $C_{\rm can}=\{c\ge0:Mc\in P\}$. Then
\begin{equation}
                     P=\sum_{g\in\{\pm1\}^3}M_gC_{\rm can},    \label{eq:polar-sum}
\end{equation}
where the right side is a Minkowski sum of eight convex cones.
Every $f\in P$ belongs to one summand if it is negative at a PR point;
otherwise it is a sum of elements from two opposite summands.
\end{proposition}
\begin{proof}
The nonsignalling polytope has the sixteen deterministic vertices and
the eight PR vertices. One short proof uses its sixteen probability
inequalities. If all marginals are interior, each $2\times2$ outcome
table has at most two zeros. A vertex needs eight independent active
constraints, so every table has two opposite zeros and gives a perfect
correlation or anticorrelation. An odd cycle parity forces all marginals
to vanish, giving a PR point; even parity leaves a nontrivial segment
of local mixtures. If any marginal is deterministic, sampling the other
party's outcomes conditionally on the remaining local variable gives
a local model, so such a vertex is deterministic.

Two distinct PR points have a local midpoint. If their patterns are
opposite, this is the uniform local correlation. Otherwise they agree
on two edges of the measurement square; the uniform distribution over
the four deterministic assignments satisfying those two edge relations
has the required midpoint. Therefore a locally nonnegative $f$ can be
negative at at most one PR point.

If $f\in P$ is negative at $R_gr$, then $R_gf$ is valid and negative
at $r$. Formula~\eqref{eq:inverse} gives $d>0$ and $c_i\ge0$, because
every deterministic neighbour lies in $\mathcal Q_{222}$. Thus
$f=M_gc$ with $c\in C_{\rm can}$.
If $f$ is nonnegative at every PR point, it is nonnegative on the
whole nonsignalling polytope. Polyhedral cone duality writes it as
a nonnegative combination of the sixteen event probabilities.
The losing events at $r$ and at $-r$ partition those sixteen events.
Each event generator is therefore the image of a source with $d=0$
in one of those two charts. Such sources are in $C_{\rm can}$.
This proves the reverse inclusion in~\eqref{eq:polar-sum}; the forward
inclusion follows from quantum validity and physical relabeling.
\end{proof}

An individual canonical cone does not contain all valid slacks: every
$M_gc$ with $c\ge0$ takes value $-d\le0$ at its own PR point,
whereas the constant slack $1$ is positive at all eight. The sum in
\eqref{eq:polar-sum} is essential.

\begin{theorem}[Exactness of the eight-orientation assembly]\label{prop:no-enlargement}
Let $K\subseteq C_{\rm can}$ be a closed convex cone containing every
nonnegative source with $d=0$. Put $P_K=\sum_gM_gK$ and
$Q_K=\{q:(1,q)\in P_K^*\}$. Then $P_K$ is closed and
\[
                         P_K\cap M(\R_+^9)=MK.
\]
Moreover
\[
 \mathcal Q_{222}\subseteq Q_K\subseteq\mathrm{NS}_{222},\qquad
 Q_K=\mathcal Q_{222}\quad\Longleftrightarrow\quad K=C_{\rm can}.
\]
\end{theorem}
\begin{proof}
The set inclusions follow from $P_K\subseteq P$ and the presence of
all sixteen event generators in $P_K$.
If $Mc\in P_K$ with $c\ge0$ and $d=0$, then $c\in K$ by assumption.
For $d>0$, write $Mc=Mc^0+\sum_{g\ne0}f_g$, where
$c^0\in K$ and $f_g=M_gc^g$.
At the own PR point $r_g$, $f_g(r_g)=-d_g\le0$. Local midpoint
positivity gives $f_g(r)\ge d_g\ge0$. In the base source coordinates,
each $f_g$ therefore has $d$-coordinate at most zero and event
coordinates $f_g(v_i)/4\ge0$. Write $c=(d,a)$ and $c^0=(d_0,a_0)$.
Then $0<d\le d_0$ and $a\ge a_0$. Thus, with $\theta=d/d_0$,
\[
               c=\theta c^0+(0,a-\theta a_0)\in K.
\]
The reverse implication is immediate.
For a convergent sequence of sums in $P_K$, their constant coefficients
are $\sum_g\sum_kc_k^g$. Nonnegativity bounds each source sequence;
closedness of $K$ gives convergent subsequences and a sum in $P_K$.
If $K=C_{\rm can}$, Proposition~\ref{prop:polar-cover} gives $Q_K=\mathcal Q_{222}$.
Otherwise every missing valid source has $d>0$ and so gives a valid
slack outside the closed cone $P_K$. A strict separating vector $z$
is in $P_K^*$. Since $P_K$ contains all probability generators,
$z$ is in the homogenized nonsignalling cone; it is nonzero and has
$z_0>0$. Normalizing $z_0=1$ gives a point outside $\mathcal Q_{222}$.
\end{proof}

\paragraph{Exact dualization.}
Let a homogeneous canonical SDP cone be given by
\[
 \mathcal L(c,y)=\sum_{k=0}^8c_k A_k+\sum_{j=1}^{p}y_j B_j,
 \qquad K=\{c:\exists y,\ \mathcal L(c,y)\succeq0\},
\]
where $y\in\R^p$. Suppose its total block order is $n$ and
$\mathcal L(c_*,y_*)\succeq\epsilon I_n$ for some $\epsilon>0$.
Then
\begin{equation}
 K^*=\{\mathcal A^*Y:\ Y\succeq0,\ \mathcal B^*Y=0\},\qquad
 (\mathcal A^*Y)_k=\operatorname{Tr}(A_kY),\quad
 (\mathcal B^*Y)_j=\operatorname{Tr}(B_jY),              \label{eq:dual}
\end{equation}
with $Y$ block diagonal in the same orders as $\mathcal L$ and no closure required.
Indeed, a convergent sequence of full adjoint images
$(\mathcal A^*Y_i,\mathcal B^*Y_i)$ satisfies
\[
 \epsilon\operatorname{Tr}Y_i
 \le\langle Y_i,\mathcal L(c_*,y_*)\rangle
 =c_*\cdot\mathcal A^*Y_i+y_*\cdot\mathcal B^*Y_i.
\]
This bounds the positive semidefinite $Y_i$ and gives a convergent
subsequence. Thus the full adjoint image is closed. Taking the dual
of the preimage cone and then the section $\mathcal B^*Y=0$ proves
\eqref{eq:dual}.

\begin{theorem}[Assembly from an exact canonical cone]\label{thm:conditional}
Suppose $K=C_{\rm can}$ and the homogeneous pencil for $K$ has a strict
feasible point. Then the full quantum set is exactly
\begin{equation}
 \mathcal Q_{222}=\{q:\ \exists Y_g\succeq0\text{ for every }g,\quad
       \mathcal B^*Y_g=0,\quad\mathcal A^*Y_g=M_g^T(1,q)\}.
                                                               \label{eq:generic-primal}
\end{equation}
It uses eight copies of the canonical dual blocks, of total order $8n$.
If the canonical pencil is rational, so is this representation.
\end{theorem}
\begin{proof}
By~\eqref{eq:dual}, the constraint for chart $g$ is exactly
$M_g^Tz\in C_{\rm can}^*$. All eight constraints are therefore equivalent
to $z\in(\sum_gM_gC_{\rm can})^*=P^*$ by
Proposition~\ref{prop:polar-cover}. Since $\mathcal Q_{222}$ is compact and
convex, $P^*\cap\{z_0=1\}=\{(1,q):q\in\mathcal Q_{222}\}$.
The eight Gram copies and their linear adjoint equations give the
claimed representation. These equations can be eliminated by fixed
linear algebra. There are no residual affine restrictions on $q$:
the homogenizations of the full-dimensional quantum set span $\R^9$,
so the prescribed adjoint targets are linearly consistent for every $q$.
\end{proof}

\paragraph{The order-$33$ specialization is an outer approximation.}
For the evaluated half-line pencil $C_{33}$, the block orders are
$(8,8,4,4,1,1,1,1,1,1,1,1,1)$, the last nine blocks being source signs.
Its supplied strict feasible point has
$c_*=(1/100,1,1,1,1,1,1,1,1)$ and
$\mathcal L(c_*,y_*)\succeq I_{33}/100$.
There are $101$ independent dual Gram entries: $92$ in the four core
blocks and nine scalar entries. The full adjoint rank is $53$.
In upper-triangular core coordinates $u\in\R^{92}$, write
$\mathcal B^*Y=B_*u$ and the core source adjoint as $A_*u$.
The rational matrix $B_*$ has rank $44$. The supplied rational matrix
$N\in\mathbb Q^{92\times48}$ is a basis of $\ker B_*$.
If $\mathcal Y(Nt)$ inserts its coordinates into symmetric blocks
of orders $8,8,4,4$, all adjoint equations are eliminated by
\begin{equation}
 \mathcal Y(Nt_g)\succeq0,\qquad
 M_g^T(1,q)-A_*Nt_g\ge0\quad\text{for all eight }g.       \label{eq:primal}
\end{equation}
The Frobenius adjoints use twice an off-diagonal upper-triangular
coefficient. Thus \eqref{eq:primal} defines an explicit outer
approximation $\widehat Q$ of total order $264$, with $104$ blocks of
maximum order $8$ and $384$ free auxiliary scalars, without coefficient
equations. It obeys
\[
                  \mathcal Q_{222}\subseteq\widehat Q
                                      \subseteq\mathrm{NS}_{222}.
\]
The first inclusion follows from canonical soundness
$C_{33}\subseteq C_{\rm can}$. For the second, $C_{33}$ contains every
nonnegative source with $d=0$: $J=K$ gives the positive event completions.
Two opposite charts therefore contain all sixteen probability generators.
Equality with $\mathcal Q_{222}$ would follow from
$C_{33}=C_{\rm can}$; no such adequacy statement is established here.
The cone $C_{33}$ is closed: bounded source coordinates bound the
diagonals and hence every entry of the two PSD $8\times8$ blocks,
so $J_0,J_2,B_G,B_H$ are bounded. The diagonal of $T_G$ then bounds
its entries, and hence $J_1$. Taking subsequences proves closedness.
Proposition~\ref{prop:no-enlargement} therefore shows that any exact
valid source excluded from $C_{33}$ makes this full outer approximation
strict; the other seven orientations cannot restore its missing inequality.

The file \texttt{full\_polar\_maps.json} gives $M,M^{-1}$, the eight
maps $M_g$, all vertex evaluations and PR-midpoint decompositions,
and the sparse rational matrices $B_*,A_*,N,A_*N$.
The accompanying exact checker reconstructs them from the event
definitions and the order-$33$ COO data. The geometric assembly theorem
applies to a different exact canonical pencil without changing any
physical relabeling or source map.
\end{document}
